\documentclass[11pt,a4paper]{article}

\usepackage[T1]{fontenc}
\usepackage[utf8]{inputenc}
\usepackage{lmodern}
\usepackage[a4paper,margin=1in]{geometry}
\usepackage{amsmath,amssymb}
\usepackage{booktabs}
\usepackage{graphicx}
\usepackage{subcaption}
\usepackage{float}
\usepackage{hyperref}
\usepackage{enumitem}

\hypersetup{
  colorlinks=true,
  linkcolor=blue,
  citecolor=blue,
  urlcolor=blue
}

\title{Hit-Competition Study at High Track Density\\\large Derived from \texttt{hit\_competition\_study.ipynb}}
\author{Quantum Track Reconstruction}
\date{\today}

% Existing figure roots from completed runs.
\newcommand{\runfigdirA}{hit_competition_outputs/20260225_115618/figures}
\newcommand{\runfigdirB}{hit_competition_outputs/20260224_140806/figures}

% Include a figure only if the file exists.
\newcommand{\safeincludegraphics}[2][]{%
  \IfFileExists{#2}{\includegraphics[#1]{#2}}{%
    \fbox{\parbox[c][0.23\textheight][c]{0.9\linewidth}{\centering Missing figure:\\ \texttt{\detokenize{#2}}}}%
  }%
}

\begin{document}
\maketitle

\begin{abstract}
This document summarizes the notebook
\texttt{Toy\_Characterisation/hit\_competition\_study.ipynb}, which tests
whether the Hamiltonian coupling term $\gamma$ suppresses false competing
segments as track density increases. The study proceeds through six stages:
(1) hit occupancy, (2) activation spectrum separation, (3) per-hit
competition behavior, (4) track-level reconstruction quality, (5)
threshold-sensitivity curves, and (6) segment-level scattering sweeps. The
workflow is designed for high-density stress testing and saves all artifacts by
run timestamp under \texttt{hit\_competition\_outputs/}.
\end{abstract}

\section{Goal and Hypothesis}
The classical solve of the linear system $Ax=b$ yields segment activations
$x_i \in [0,1]$. Segments sharing a hit are coupled through the Hamiltonian.
The central hypothesis is:
\begin{quote}
As occupancy increases, true segments remain near or above the operating
threshold, while false competitors are suppressed by the competition structure
encoded in $A$.
\end{quote}

The notebook evaluates this mechanism at multiple event densities and angular
phase-space settings.

\section{Configuration and Definitions}
\subsection{Detector and Event Generation}
The notebook uses a 5-module planar geometry with module spacing
$\Delta z = 33\,\mathrm{mm}$ and symmetric angular generation windows
$\phi_{\max}=\theta_{\max}\in\{0.2,\,0.1,\,0.04\}\,\mathrm{rad}$.
A density sweep is performed over
\[
N_\text{tracks}\in\{5,10,20,30,50,75,100,150\}
\]
with repeated event generation per point.

\subsection{Hamiltonian Parameters}
The baseline parameter set is
\[
\sigma_{\mathrm{res}}=0,\quad
\sigma_{\mathrm{scatt}}=10^{-4}\,\mathrm{rad},\quad
\gamma=1.5,\quad
\delta=1.0,
\]
and
\[
\varepsilon=\sqrt{2\theta_s^2 + 12\theta_r^2 + 2\theta_{\min}^2},
\quad
\theta_s=s\,\sigma_{\mathrm{scatt}},
\quad
\theta_r=\arctan\!\left(\frac{s\,\sigma_{\mathrm{res}}}{\Delta z}\right),
\]
with scale $s=3$. The activation operating point is
\[
\text{baseline}=\frac{\delta}{\delta+\gamma},
\qquad
\text{threshold}=\frac{1+\text{baseline}}{2}.
\]

\section{Analysis Workflow}
\subsection{Step 1: Hit Occupancy}
For each generated event, the study builds a map
\texttt{hit\_id} $\to$ candidate segment indices and records occupancy
statistics (mean, median, max, spread). This quantifies combinatorial pressure
and validates the expected scaling for internal layers.

\subsection{Step 2: Activation Spectrum (True vs False)}
Each candidate segment is labeled true/false from ground-truth track adjacency.
The solved activations $x_i$ are then compared between classes, including:
\begin{itemize}[leftmargin=1.5em]
  \item fraction of true segments above threshold,
  \item fraction of false segments above threshold,
  \item separation statistic from pooled variance.
\end{itemize}
This step directly measures whether ambiguity suppression degrades at high
occupancy.

\subsection{Step 3: Per-Hit Competition Curves}
For contested hits (shared by multiple segments), the notebook compares:
\begin{itemize}[leftmargin=1.5em]
  \item true segment activation $x_{\text{true}}$,
  \item summed false competitor load $\sum_j x_{\text{false},j}$,
  \item number of competitors.
\end{itemize}
Scatter and summary curves test whether true segments remain robust under
increasing local competition.

\subsection{Step 4: Track-Level Reconstruction Quality}
An end-to-end chain
\texttt{generate $\to$ construct Hamiltonian $\to$ solve $\to$ get\_tracks $\to$ validate}
is run for each density and angle. Metrics include:
\begin{itemize}[leftmargin=1.5em]
  \item reconstruction efficiency,
  \item ghost rate,
  \item clone fraction.
\end{itemize}
Truth and reconstructed events are also exported as JSON snapshots.

\subsection{Step 5: Threshold Sensitivity (ROC-like)}
Thresholds are swept over
\[
\tau\in[0.2,0.9]
\]
at representative densities. For each $\tau$, efficiency and ghost rate are
computed, giving operating curves and highlighting the default threshold point.

\subsection{Step 6: Segment-Level Scattering Sweep}
A second scan varies multiple-scattering strength by factors
$1\times$, $2\times$, and $4\times$ (with corresponding $\varepsilon$ values),
across event sizes 10 to 100 tracks. Reported metrics:
\begin{itemize}[leftmargin=1.5em]
  \item segment efficiency,
  \item segment false rate,
  \item true/false segment-pair counts,
  \item accepted true/false pairs with $\theta\le\varepsilon$.
\end{itemize}
Additional histograms compare true and false pairwise-angle distributions for
fixed-size high-density ensembles.

\section{Produced Artifacts}
Each run creates a timestamped directory:
\begin{center}
\texttt{hit\_competition\_outputs/<RUN\_TAG>/\{data,figures,events\}}
\end{center}
Typical figures (saved when corresponding steps are executed):
\begin{itemize}[leftmargin=1.5em]
  \item \texttt{step1\_occupancy\_comparison.png}
  \item \texttt{step1\_occupancy\_histograms.png}
  \item \texttt{step2\_activation\_spectra.png}
  \item \texttt{step2\_activation\_summary.png}
  \item \texttt{step3\_competition\_scatter.png}
  \item \texttt{step3\_competition\_summary.png}
  \item \texttt{step4\_reco\_quality.png}
  \item \texttt{step5\_roc\_by\_angle.png}
  \item \texttt{step5\_threshold\_curves.png}
  \item \texttt{step6\_scattering\_comparison.png}
  \item \texttt{step6\_acceptance\_histograms.png}
\end{itemize}

\section{Representative Figures}
\begin{figure}[H]
  \centering
  \begin{subfigure}[b]{0.48\textwidth}
    \centering
    \safeincludegraphics[width=\linewidth]{\runfigdirA/step1_occupancy_comparison.png}
    \caption{Step 1 occupancy comparison.}
  \end{subfigure}\hfill
  \begin{subfigure}[b]{0.48\textwidth}
    \centering
    \safeincludegraphics[width=\linewidth]{\runfigdirA/step2_activation_summary.png}
    \caption{Step 2 activation separation summary.}
  \end{subfigure}
  \caption{Occupancy and activation-separation diagnostics.}
\end{figure}

\begin{figure}[H]
  \centering
  \begin{subfigure}[b]{0.48\textwidth}
    \centering
    \safeincludegraphics[width=\linewidth]{\runfigdirA/step3_competition_summary.png}
    \caption{Step 3 per-hit competition summary.}
  \end{subfigure}\hfill
  \begin{subfigure}[b]{0.48\textwidth}
    \centering
    \safeincludegraphics[width=\linewidth]{\runfigdirB/step4_reco_quality.png}
    \caption{Step 4 track-level quality.}
  \end{subfigure}
  \caption{Competition behavior and end-to-end reconstruction quality.}
\end{figure}

\begin{figure}[H]
  \centering
  \begin{subfigure}[b]{0.48\textwidth}
    \centering
    \safeincludegraphics[width=\linewidth]{seg_eff_1mrad.png}
    \caption{Segment efficiency reference plot.}
  \end{subfigure}\hfill
  \begin{subfigure}[b]{0.48\textwidth}
    \centering
    \safeincludegraphics[width=\linewidth]{seg_hist_1mrad.png}
    \caption{Segment-angle histogram reference plot.}
  \end{subfigure}
  \caption{Additional segment-level reference outputs available in the study directory.}
\end{figure}

\section{Interpretation Guide}
The notebook supports the following checks:
\begin{enumerate}[leftmargin=1.5em]
  \item \textbf{Combinatorics scaling}: occupancy trends with density and angle
  indicate where overlap becomes severe.
  \item \textbf{Activation robustness}: true/false spectrum separation indicates
  whether competition suppression remains effective.
  \item \textbf{Operational envelope}: ROC-like curves identify threshold headroom
  as density increases.
  \item \textbf{Physics stress response}: scattering sweeps quantify sensitivity of
  segment-level acceptance and contamination.
\end{enumerate}

\section{Reproducibility Notes}
\begin{itemize}[leftmargin=1.5em]
  \item The notebook can either compute fresh results or load a previous run via
  \texttt{LOAD\_RUN}.
  \item All JSON summaries are saved in the run's \texttt{data/} directory.
  \item To display figures from another run, update
  \texttt{\textbackslash runfigdirA} and \texttt{\textbackslash runfigdirB} at the top of this file.
\end{itemize}

\section{Conclusion}
This study provides a structured, multi-level validation of hit-competition
handling in the Hamiltonian reconstruction approach. By combining local
(occupancy and per-hit competition), intermediate (segment activations), and
global (track-level metrics and threshold scans) diagnostics across density,
angle phase space, and scattering strength, the notebook establishes a clear
framework for quantifying high-density robustness.

\end{document}
